% ===================================================================
%  ਸਾਰਣੀ ਨੰ. 15 — ਮੰਡੀ। — ਖਾਣ ਦੀਆਂ ਚੀਜ਼ਾਂ।
%  Traduction 2026 d'apres texto/io, controlee sur texto/fr, texto/en
%  et texto/hi. Consigne : voir texto/pa/10-sarni-01.tex.
% ===================================================================

%%K t15-apar-1 apar
\VUsaut{4.0mm}
\VUtitre{15.0pt}{60}{ਸਾਰਣੀ ਨੰ. 15}
\VUsaut{3.0mm}

%%K t15-tit-1 sub
\VUpk{11.4pt}{40}{ਮੰਡੀ। --- ਖਾਣ ਦੀਆਂ ਚੀਜ਼ਾਂ।}
\VUsaut{3.0mm}

%%K t15-01-1 p
1. --- ਜੋ ਵਪਾਰੀ ਸਾਨੂੰ ਸਾਡੀ ਲੋੜ ਦਾ ਖਾਣਾ ਦਿੰਦੇ ਹਨ, ਉਨ੍ਹਾਂ ਵਿਚੋਂ ਬਹੁਤੇ
\VUgras{ਢਕੀ ਮੰਡੀ} ਦੇ \VUgras{ਹਾਲ} \textsuperscript{(1)} ਦੇ ਹੇਠਾਂ ਜਮ੍ਹਾਂ ਰਹਿੰਦੇ
ਹਨ, ਜਾਂ ਕੋਲ ਦੀਆਂ ਗਲੀਆਂ ਵਿਚ।

%%K t15-02-1 p
2. --- \VUgras{ਸੂਰ ਦੇ ਮਾਸ ਦਾ ਕਸਾਈ} \textsuperscript{(2)}, ਜੋ \VUgras{ਹੈਮ}
\textsuperscript{(3)}, \VUgras{ਲਹੂ ਦੀ ਸੌਸੇਜ} \textsuperscript{(4)},
\VUgras{ਸੌਸੇਜ} \textsuperscript{(5)}, \VUgras{ਆਂਦਰਾਂ ਦੀ ਸੌਸੇਜ}
\textsuperscript{(6)} ਤੇ \VUgras{ਚਰਬੀ} \textsuperscript{(7)} ਵੇਚਦਾ ਹੈ, ਉਹ
\VUgras{ਮੱਖਣ ਤੇ ਪਨੀਰ ਵੇਚਣ ਵਾਲੀ} \textsuperscript{(8)} ਕੋਲ ਬਹਿੰਦਾ ਹੈ, ਜਿਸ ਦਾ ਠੇਲਾ ਲਾਲ
ਛਿੱਲ ਵਾਲੇ \VUgras{ਹਾਲੈਂਡ ਦੇ ਪਨੀਰਾਂ} \textsuperscript{(9)},
\VUgras{ਕਾਮਾਂਬੇਰ} \textsuperscript{(10)}, \VUgras{ਗਰੂਯੇਰ ਦੇ ਵੱਡੇ ਚੱਕਰਾਂ}
\textsuperscript{(11)} ਤੇ ਤਾਜ਼ੇ \VUgras{ਮੱਖਣ ਦੀਆਂ ਪੇੜੀਆਂ} \textsuperscript{(12)} ਨਾਲ
ਸਜਿਆ ਹੈ।

%%K t15-03-1 p
3. --- ਉਸ ਦੇ ਸਾਹਮਣੇ ਇਕ \VUgras{ਸਬਜ਼ੀ ਵੇਚਣ ਵਾਲੀ} \textsuperscript{(13)} ਆਪਣੀਆਂ
ਗਾਹਕਾਂ ਨੂੰ ਸੋਹਣੇ \VUgras{ਟਮਾਟਰ} \textsuperscript{(14)}, \VUgras{ਫੁੱਲ ਗੋਭੀ}
\textsuperscript{(15)}, \VUgras{ਸਲਾਦ} \textsuperscript{(16)}, \VUgras{ਬੰਦ ਗੋਭੀ}
\textsuperscript{(17)}, \VUgras{ਕੱਦੂ} \textsuperscript{(19)}, \VUgras{ਖ਼ਰਬੂਜ਼ੇ}
\textsuperscript{(18)} ਤੇ \VUgras{ਖੁੰਬਾਂ} \textsuperscript{(20)} ਤਕ ਵਿਖਾਉਂਦੀ ਹੈ। ਪਰ
ਇਕ \VUgras{ਸ਼ਰਾਬ ਵਾਲੇ} ਨੇ ਆਪਣੀ \VUgras{ਢੋਣ ਦੀ ਗੱਡੀ} \textsuperscript{(21)}, ਜੋ
\VUgras{ਬੀਅਰ ਦੇ ਪੀਪਿਆਂ} \textsuperscript{(22)} ਨਾਲ ਲੱਦੀ ਹੈ, ਠੀਕ ਉਸ ਦੇ ਠੇਲੇ ਦੇ ਸਾਹਮਣੇ
ਖੜੀ ਕਰ ਦਿੱਤੀ ਹੈ ਤੇ ਉਸ ਦੀਆਂ ਗਾਹਕਾਂ ਨੂੰ ਕੋਲ ਨਹੀਂ ਆਉਣ ਦਿੰਦੀ। ਇਕ ਬਾਗ਼ਬਾਨ ਉਸ ਨੂੰ
\VUgras{ਹਾਥੀ-ਚੱਕ} \textsuperscript{(23)} ਦੀ ਇਕ ਟੋਕਰੀ ਲਿਆ ਰਿਹਾ ਹੈ।

%%K t15-tit-2 sub
\VUsaut{4.0mm}
\VUpk{10.2pt}{40}{ਵੇਚਣਾ ਤੇ ਖ਼ਰੀਦਣਾ।}
\VUsaut{3.0mm}

%%K t15-04-1 p
4. --- ਮੰਡੀ ਵਿਚ ਵੱਡੀ ਚਹਿਲ-ਪਹਿਲ ਹੈ, ਕਿਉਂਕਿ \VUgras{ਨੀਲਾਮੀ} \textsuperscript{(24)}
ਦਾ ਵੇਲਾ ਆ ਗਿਆ ਹੈ। ਜੋ \VUgras{ਘਰ ਵਾਲੀਆਂ} \textsuperscript{(26)} ਇੱਥੇ ਸੌਦਾ ਕਰਨ
ਆਉਂਦੀਆਂ ਹਨ, ਉਹ ਰਾਹ ਵਿਚ \VUgras{ਓਝੜੀ ਵੇਚਣ ਵਾਲੀ} \textsuperscript{(25)} ਦੇ ਠੇਲੇ ਦੇ
ਸਾਹਮਣੇ ਰੁਕ ਕੇ \VUgras{ਓਝੜੀ} ਜਾਂ ਇਕ \VUgras{ਵੱਛੇ ਦਾ ਸਿਰ} ਖ਼ਰੀਦਦੀਆਂ ਹਨ।

%%K t15-05-1 p
5. --- ਇਕ ਮੋਟਾ \VUgras{ਰਸੋਈਆ} \textsuperscript{(27)} \VUgras{ਮੱਛੀ ਵੇਚਣ ਵਾਲੀ}
\textsuperscript{(28)} ਨਾਲ ਇਕ ਬਹੁਤ ਸੋਹਣੀ \VUgras{ਟੂਨਾ} ਦਾ ਮੁੱਲ-ਭਾਅ ਕਰ ਰਿਹਾ ਹੈ, ਜੋ
ਆਪਣੇ ਮਨ ਨਾਲ ਭਾਅ ਵਧਾਉਂਦੀ ਜਾਂਦੀ ਹੈ। ਉਸ ਕੋਲ \VUgras{ਬਾਮ, ਸੋਲ, ਟ੍ਰਾਊਟ} ਤੇ
\VUgras{ਕਾਰਪ} ਦੀ ਵੱਡੀ ਖੇਪ ਆਈ ਹੈ। ਉਸ ਦਾ \VUgras{ਕਾਡ} ਤੇ ਉਸ ਦੀਆਂ \VUgras{ਸਿੱਪੀਆਂ}
ਬਹੁਤ ਤਾਜ਼ੀਆਂ ਹਨ।

%%K t15-tit-3 sub
\VUsaut{4.0mm}
\VUpk{10.2pt}{40}{ਸਸਤਾ ਮਾਲ ਤੇ ਮਹਿੰਗਾ।}
\VUsaut{3.0mm}

%%K t15-06-1 p
6. --- ਉਸ ਕੋਲ ਬੈਠੀ \VUgras{ਕੁਕੜੀਆਂ ਵੇਚਣ ਵਾਲੀ} \textsuperscript{(29)} ਬਹੁਤ ਇਮਾਨਦਾਰ
ਹੈ। ਜਦੋਂ ਉਹ ਕੋਈ \VUgras{ਟਰਕੀ}, ਕੋਈ \VUgras{ਹੰਸ}, ਕੋਈ \VUgras{ਬੱਤਖ} ਜਾਂ ਸਧਾਰਨ
\VUgras{ਕੁਕੜੀ} ਵੇਚਦੀ ਹੈ, ਤਾਂ ਵਾਜਬ ਭਾਅ ਤੋਂ ਵੱਧ ਨਹੀਂ ਮੰਗਦੀ।

%%K t15-07-1 p
7. --- ਚੌਕ ਦੇ ਕੋਨੇ ਉੱਤੇ, ਪਹਿਲੀ ਮੰਜ਼ਿਲ ਉੱਤੇ, ਇਕ \VUgras{ਕੱਪੜੇ ਦੀ ਦੁਕਾਨ ਵਾਲੀ}
\textsuperscript{(30)} ਨੇ ਹਾਲ ਹੀ ਦੁਕਾਨ ਖੋਲ੍ਹੀ ਹੈ, ਜਿੱਥੇ ਖ਼ਰੀਦਣ ਵਾਲੀਆਂ ਨੂੰ
\VUgras{ਮੇਜ਼ਪੋਸ਼}, \VUgras{ਰੁਮਾਲ, ਪਰਨੇ} ਤੇ \VUgras{ਤੌਲੀਆਂ} ਦੀ ਬਹੁਤ ਵਧੀਆ ਚੋਣ
ਹਮੇਸ਼ਾ ਮਿਲਦੀ ਹੈ। ਉਸੇ ਮੰਜ਼ਿਲ ਉੱਤੇ ਇਕ \VUgras{ਛੁਰੀ ਬਣਾਉਣ ਵਾਲੇ} \textsuperscript{(31)}
ਨੇ ਆਪਣੀ ਕਾਰਜਸ਼ਾਲਾ ਲਾਈ ਹੈ। ਉਹ ਮੰਡੀ ਦੀਆਂ ਤੀਵੀਂਆਂ ਦੀਆਂ \VUgras{ਛੁਰੀਆਂ} ਤੇਜ਼ ਕਰਦਾ ਹੈ
ਤੇ ਮਾਲੀ ਲਈ ਸਭ ਤੋਂ ਕਰੜੇ \VUgras{ਫ਼ੌਲਾਦ} ਤੋਂ \VUgras{ਛਾਂਗਣ ਦੀਆਂ ਦਾਤਰੀਆਂ} ਬਣਾਉਂਦਾ
ਹੈ।

%%K t15-tit-4 sub
\VUsaut{4.0mm}
\VUpk{10.2pt}{40}{ਕਿਰਿਆਨੇ ਦੀ ਦੁਕਾਨ।}
\VUsaut{3.0mm}

%%K t15-08-1 p
8. --- ਉਸ ਦੀ ਕਾਰਜਸ਼ਾਲਾ ਦੇ ਹੇਠਾਂ ਕੁਝ ਵੇਲਾ ਪਹਿਲਾਂ ਇਕ ਵੱਡੀ ਖਾਣੇ ਦੀ ਦੁਕਾਨ ਖੁੱਲ੍ਹੀ ਹੈ।
ਇਹ ਹੁਣ ਪੁਰਾਣੇ ਦਿਨਾਂ ਵਾਲੀ ਉਹ ਸਾਦੀ ਨਿੱਕੀ \VUgras{ਕਿਰਿਆਨੇ ਦੀ ਦੁਕਾਨ}
\textsuperscript{(32)} ਨਹੀਂ ਰਹੀ, ਜੋ ਸਿਰਫ਼ ਰੋਜ਼ ਦਾ ਮਾਲ ਵੇਚਦੀ ਸੀ — \VUgras{ਚਾਹ}
\textsuperscript{(33)}, \VUgras{ਚਾਕਲੇਟ} \textsuperscript{(34)},
\VUgras{ਡਲੀ ਦੀ ਖੰਡ} \textsuperscript{(37)} ਤੇ \VUgras{ਸਾਬਣ} \textsuperscript{(38)}।
ਇੱਥੇ ਸਭ ਕੁਝ ਵਿਕਦਾ ਹੈ : \VUgras{ਕੁਰਕੁਰੇ ਬਿਸਕੁਟ}, \VUgras{ਡੱਬਾਬੰਦ ਚੀਜ਼ਾਂ}
\textsuperscript{(35)}, \VUgras{ਸ਼ਰਾਬਾਂ} \textsuperscript{(36)}, \VUgras{ਰਮ, ਬਰਾਂਡੀ,
ਕਿਰਸ਼, ਅਨੀਸੇਤ} ਤੇ \VUgras{ਕੁਰਾਸਾਓ}। ਇੱਥੇ ਸ਼ਿਕਾਰ ਵੀ ਓਨੀ ਹੀ ਸੌਖ ਨਾਲ ਮਿਲਦਾ ਹੈ —
\VUgras{ਸਹਾ} \textsuperscript{(39)}, \VUgras{ਚਿੜੀਆਂ} \textsuperscript{(40)},
\VUgras{ਤਿੱਤਰ} \textsuperscript{(41)}, \VUgras{ਬਟੇਰ} \textsuperscript{(42)},
\VUgras{ਜੰਗਲੀ ਬੱਤਖ} \textsuperscript{(43)} ਜਾਂ \VUgras{ਤਿੱਤਰ-ਮੋਰ}
\textsuperscript{(44)} — ਜਿੰਨੇ ਨਿੱਘੇ ਦੇਸਾਂ ਦੇ ਫਲ, ਜਿਵੇਂ \VUgras{ਕੇਲੇ}
\textsuperscript{(46)} ਤੇ \VUgras{ਅਨਾਨਾਸ}।

%%K t15-09-1 p
9. --- ਇਕ ਬਹੁਤ ਉਪਕਾਰੀ ਕਿਰਿਆਨੇ ਦਾ \VUgras{ਕਾਰਿੰਦਾ} \textsuperscript{(45)} ਜੋ ਮੰਗਿਆ
ਜਾਵੇ ਉਹੋ ਦੇਣ ਨੂੰ ਤਿਆਰ ਰਹਿੰਦਾ ਹੈ : \VUgras{ਮੁਰੱਬਾ} \textsuperscript{(47)}, ਕਰੌਂਦੇ
ਦਾ ਜਾਂ ਬਿਹੀ ਦਾ, \VUgras{ਚਰਬੀ} \textsuperscript{(48)}, \VUgras{ਖੀਰੇ}
\textsuperscript{(49)} ਜਾਂ ਲੂਣੀ \VUgras{ਸਾਰਡੀਨ} \textsuperscript{(50)}।

%%K t15-09-2 p
ਇਹ \VUgras{ਗਾਹਕਾਂ} \textsuperscript{(51)} ਲਈ ਬਹੁਤ ਸੌਖ ਦੀ ਗੱਲ ਹੈ, ਜੋ ਸਭ ਤੋਂ ਸਧਾਰਨ ਮਾਲ
ਦੇ ਕੋਲ ਹੀ ਸਭ ਤੋਂ ਟਾਵਾਂ ਅਗੇਤਾ ਮਾਲ ਤੇ ਸਭ ਤੋਂ ਸੁਆਦੀ ਫਲ ਲੱਭਦੀਆਂ ਹਨ : ਰਸੀਲੀਆਂ
\VUgras{ਨਾਸ਼ਪਾਤੀਆਂ} \textsuperscript{(52)}, \VUgras{ਆਲੂਬੁਖ਼ਾਰੇ}
\textsuperscript{(53)}, \VUgras{ਅੰਗੂਰ} \textsuperscript{(54)}, \VUgras{ਆੜੂ}
\textsuperscript{(55)}, \VUgras{ਬਦਾਮ} \textsuperscript{(56)} ਤੇ \VUgras{ਅਖ਼ਰੋਟ}
\textsuperscript{(57)}।

%%K t15-tit-5 sub
\VUsaut{4.0mm}
\VUpk{10.2pt}{40}{ਗਾਹਕ।}
\VUsaut{3.0mm}

%%K t15-10-1 p
10. --- \VUgras{ਚੌਲ} \textsuperscript{(58)} ਤੇ \VUgras{ਮਸਰ} \textsuperscript{(59)}
ਧੂੰਏਂ ਦਿੱਤੀ \VUgras{ਹੈਰਿੰਗ} \textsuperscript{(60)} ਦੇ ਡੱਬੇ ਕੋਲ ਰੱਖੇ ਹਨ;
\VUgras{ਆਲੂ} \textsuperscript{(61)} ਤੇ \VUgras{ਸੇਮ} \textsuperscript{(63)}
\VUgras{ਸੇਬਾਂ} \textsuperscript{(62)}, \VUgras{ਅੰਜੀਰਾਂ} \textsuperscript{(64)},
\VUgras{ਸੁੱਕੇ ਆਲੂਬੁਖ਼ਾਰਿਆਂ} \textsuperscript{(65)} ਤੇ \VUgras{ਪਹਾੜੀ ਬਦਾਮਾਂ}
\textsuperscript{(66)} ਕੋਲ ਬੈਠੇ ਹਨ। ਇਸ ਦੁਕਾਨ ਵਿਚ ਤਾਜ਼ੇ \VUgras{ਆਂਡੇ}
\textsuperscript{(67)} ਤੇ \VUgras{ਜ਼ੈਤੂਨ} \textsuperscript{(68)} ਵੀ ਮਿਲਦੇ ਹਨ, ਸੋ
\VUgras{ਟੋਕਰੀਆਂ} \textsuperscript{(70)} ਤੇ \VUgras{ਜਾਲੀ ਦੇ ਥੈਲੇ}
\textsuperscript{(73)} ਬਹੁਤ ਛੇਤੀ ਭਰ ਜਾਂਦੇ ਹਨ।

%%K t15-11-1 p
11. --- ਇਸ ਵਧੀਆ ਦੁਕਾਨ ਦੀ \VUgras{ਕੌਫ਼ੀ} \textsuperscript{(72)} ਨੇ
\VUgras{ਨੌਕਰਾਣੀਆਂ} \textsuperscript{(69)} ਤੇ ਰਸੋਇਣਾਂ ਵਿਚ ਚੰਗਾ-ਖ਼ਾਸਾ ਨਾਂ ਕਮਾਇਆ ਹੈ,
ਕਿਉਂਕਿ ਬੁੱਢਾ \VUgras{ਕਿਰਿਆਨੇ ਵਾਲਾ} \textsuperscript{(71)} ਉਸ ਨੂੰ ਆਪ ਬੜੇ ਧਿਆਨ ਨਾਲ
ਭੁੰਨਦਾ ਹੈ।

%%K t15-tit-6 sub
\VUsaut{4.0mm}
\VUpk{10.2pt}{40}{ਵੇਖਣ ਦੀਆਂ ਚੀਜ਼ਾਂ।}
\VUsaut{3.0mm}

%%K t15-12-1 p
12. --- ਸਾਰੇ ਲੋਕ ਮੰਡੀ ਸੌਦਾ ਕਰਨ ਹੀ ਨਹੀਂ ਆਉਂਦੇ। ਹਾਲ ਤਕ ਜਾਂਦੀ ਉਸ ਨਿੱਕੀ ਗਲੀ ਦੀ ਸੋਹਣੀ
ਆਵਾਜਾਈ ਵਿਚ ਕਈ ਤਰ੍ਹਾਂ ਦੇ ਲੋਕ ਵਿਖਾਈ ਦਿੰਦੇ ਹਨ। ਇਕ ਭਲੇ \VUgras{ਬੁੱਚੜਖ਼ਾਨੇ ਦੇ ਮਜ਼ਦੂਰ}
\textsuperscript{(75)} ਕੋਲ, ਜੋ ਆਪਣੇ ਅੱਗੇ ਇਕ \VUgras{ਠੇਲਾ} \textsuperscript{(74)}
ਧੱਕ ਰਿਹਾ ਹੈ, ਇਹ ਰਹੇ ਦੋ ਛੁੱਟੀ ਉੱਤੇ ਆਏ ਸਿਪਾਹੀ, ਜਿਨ੍ਹਾਂ ਨੂੰ ਆਪਣੀ ਚਮਕਦੀ ਵਰਦੀ ਵਿਖਾਉਣ
ਦਾ ਬੜਾ ਮਾਣ ਹੈ : \VUgras{ਹੁਸਾਰ} \textsuperscript{(76)} ਆਪਣੇ ਨੀਲੇ \VUgras{ਦੋਲਮਾਨ} ਤੇ
\VUgras{ਖੰਭਾਂ ਵਾਲੇ ਸ਼ਾਕੋ} ਵਿਚ, ਤੇ \VUgras{ਜ਼ੁਆਵ ਦਾ ਹੌਲਦਾਰ} ਢਿੱਲੀ ਨਿੱਕਰ ਵਿਚ, ਸਿਰ
ਉੱਤੇ \VUgras{ਫ਼ੇਜ਼} ਲਾਈ।

%%K t15-13-1 p
13. --- \VUgras{ਨਾਨਬਾਈ} \textsuperscript{(80)}, ਜੋ \VUgras{ਨਾਨਬਾਈਖ਼ਾਨੇ}
\textsuperscript{(78)} ਦੀ ਦਹਿਲੀਜ਼ ਉੱਤੇ ਖੜਾ ਹੈ, ਉਸ ਨੇ ਹੁਣੇ ਹੁਣੇ ਆਪਣੀ \VUgras{ਰੋਟੀ}
\textsuperscript{(79)} ਦਾ ਆਟਾ ਗੁੰਨ੍ਹਿਆ ਹੈ ਤੇ ਤੰਦੂਰ ਤੋਂ ਉਹ \VUgras{ਕਰੋਸਾਂ}
\textsuperscript{(81)}, \VUgras{ਨਿੱਕੀਆਂ ਰੋਟੀਆਂ} \textsuperscript{(82)} ਤੇ
\VUgras{ਗੋਲ ਰੋਟੀਆਂ} \textsuperscript{(83)} ਕੱਢੀਆਂ ਹਨ ਜੋ ਖਿੜਕੀ ਵਿਚ ਰੱਖੀਆਂ ਹਨ।

%%K t15-14-1 p
14. --- ਨਾਨਬਾਈਖ਼ਾਨੇ ਦੇ ਉੱਤੇ ਇਕ ਸਿਆਣੀ \VUgras{ਸਿਲਾਈ ਵਾਲੀ} \textsuperscript{(84)}
ਰਹਿੰਦੀ ਹੈ, ਜੋ ਕਈ \VUgras{ਕਢਾਈ ਕਰਨ ਵਾਲੀਆਂ} \textsuperscript{(85)} ਰੱਖਦੀ ਹੈ। ਉਸ ਕੋਲ
ਇਕ \VUgras{ਧੋਬਣ ਤੇ ਇਸਤਰੀ ਕਰਨ ਵਾਲੀ} \textsuperscript{(86)} ਰਹਿੰਦੀ ਹੈ, ਜੋ
\VUgras{ਕੱਪੜੇ} \textsuperscript{(88)} ਧੋ ਕੇ ਤੇ ਸੁਕਾ ਕੇ ਉਨ੍ਹਾਂ ਨੂੰ ਕਲਫ਼ ਲਾਉਂਦੀ ਹੈ ਤੇ
\VUgras{ਇਸਤਰੀ} \textsuperscript{(87)} ਨਾਲ ਦਬਦੀ ਹੈ।

%%K t15-tit-7 sub
\VUsaut{4.0mm}
\VUpk{10.2pt}{40}{ਮਾਸ, ਮੱਛੀ ਤੇ ਸਬਜ਼ੀ।}
\VUsaut{3.0mm}

%%K t15-15-1 p
15. --- ਹੇਠਲੀ ਮੰਜ਼ਿਲ ਦੀ \VUgras{ਕਸਾਈ ਦੀ ਦੁਕਾਨ} \textsuperscript{(89)} ਵਿਚ ਹਰ ਤਰ੍ਹਾਂ
ਦਾ ਮਾਸ ਵਿਕਦਾ ਹੈ — \VUgras{ਭੇਡ ਦਾ} \textsuperscript{(90)}, \VUgras{ਗਾਂ ਦਾ}
\textsuperscript{(94)} ਜਾਂ \VUgras{ਵੱਛੇ ਦਾ} \textsuperscript{(92)} — \VUgras{ਰਾਣਾਂ,
ਚਾਪਾਂ, ਟੁਕੜੇ} ਤੇ \VUgras{ਪਸਲੀਆਂ} ਦੇ ਰੂਪ ਵਿਚ। \VUgras{ਕਸਾਈ ਦਾ ਛੋਕਰਾ}
\textsuperscript{(93)} ਇਹ ਸਾਰਾ ਮਾਸ ਕੱਟਦਾ ਹੈ ਤੇ \VUgras{ਗੁੱਦੇ ਦੀਆਂ ਹੱਡੀਆਂ}
\textsuperscript{(96)} ਵੱਖ ਰੱਖਦਾ ਹੈ। ਕਸਾਈ ਦੇ ਬੂਹੇ ਉੱਤੇ ਇਕ
\VUgras{ਸਿੱਪੀਆਂ ਵੇਚਣ ਵਾਲੀ} \textsuperscript{(97)} ਖੜੀ ਹੈ, ਜੋ \VUgras{ਸਿੱਪੀਆਂ}
\textsuperscript{(98)} ਵੇਚਦੀ ਹੈ। ਚਾਹੋ ਤਾਂ ਉਹ ਉਨ੍ਹਾਂ ਨੂੰ ਖੋਲ੍ਹ ਵੀ ਦਿੰਦੀ ਹੈ।

%%K t15-16-1 p
16. --- ਇਹ ਸਾਰੇ ਵਪਾਰੀ ਲਾਇਸੰਸ ਜਾਂ ਥਾਂ ਦਾ ਮਹਿਸੂਲ ਦਿੰਦੇ ਹਨ; ਇਸੇ ਲਈ
\VUgras{ਸਿਪਾਹੀ} \textsuperscript{(100)} ਉਨ੍ਹਾਂ ਵਿਚਾਰੀਆਂ
\VUgras{ਘੁੰਮ ਕੇ ਸਬਜ਼ੀ ਵੇਚਣ ਵਾਲੀਆਂ} \textsuperscript{(99)} ਦਾ ਬਿਨਾਂ ਤਰਸ ਦੇ ਪਿੱਛਾ
ਕਰਦਾ ਹੈ, ਜੋ ਆਪਣੀਆਂ ਰੇੜ੍ਹੀਆਂ ਉੱਤੇ \VUgras{ਗਾਜਰਾਂ, ਮੂਲੀਆਂ, ਸ਼ਲਗਮ, ਗੰਦਨਾ} ਜਾਂ
\VUgras{ਸ਼ਤਾਵਰੀ ਦੀਆਂ ਗੱਠੜੀਆਂ, ਤਾਜ਼ੇ ਮਟਰ, ਪਾਲਕ, ਖੱਟੀ ਬੂਟੀ, ਜਲ-ਕੁੰਭੀ} ਤੇ
\VUgras{ਧਨੀਆ} ਢੋ ਕੇ ਉਨ੍ਹਾਂ ਨਾਲ ਮੁਕਾਬਲਾ ਕਰਦੀਆਂ ਹਨ।

%%K t15-tit-8 sub
\VUsaut{4.0mm}
\VUpk{10.2pt}{40}{ਸਿਪਾਹੀ ਤੋਂ ਬਚੋ!}
\VUsaut{3.0mm}

%%K t15-17-1 p
17. --- \VUgras{ਥੋਮ} \textsuperscript{(101)} ਤੇ \VUgras{ਗੰਢੇ} ਵੇਚਣ ਵਾਲਾ ਮੁੰਡਾ ਖ਼ੂਬ
ਜਾਣਦਾ ਹੈ ਕਿ ਉਸ ਨੂੰ ਵੀ ਮੰਡੀ ਕੋਲ ਆਪਣਾ ਮਾਲ ਵੇਚਣ ਦੀ ਆਗਿਆ ਨਹੀਂ। ਉਹ ਨੱਸ ਖੜਦਾ ਹੈ, ਇਸ
ਖ਼ਤਰੇ ਨਾਲ ਕਿ \VUgras{ਝੀਂਗਾ-ਮੱਛੀ} ਤੇ \VUgras{ਲਾਬਸਟਰ} \VUgras{ਵੇਚਣ ਵਾਲੀ}
\textsuperscript{(102)} ਦੀ \VUgras{ਕੇਕੜਿਆਂ}, \VUgras{ਕ੍ਰੇਫ਼ਿਸ਼} ਤੇ
\VUgras{ਝੀਂਗਿਆਂ} ਦੀ ਟੋਕਰੀ ਉਲਟਾ ਦੇਵੇ।

%%K t15-18-1 p
18. --- ਸਾਰੇ ਲੋਕ ਕਾਹਲੀ ਵਿਚ ਹਨ ਤੇ ਵੱਡਾ ਰੌਲਾ ਪਾ ਰਹੇ ਹਨ; ਪਰ ਇਸ ਨਾਲ ਘੁੰਮਦੇ
\VUgras{ਸ਼ੀਸ਼ਾ-ਸਾਜ਼} \textsuperscript{(103)} ਦੀ ਹਾਕ ਸਾਫ਼ ਸੁਣਾਈ ਦੇਣੋਂ ਨਹੀਂ ਰੁਕਦੀ, ਨਾ
ਉਸ \VUgras{ਕੋਲੇ ਵਾਲੇ} \textsuperscript{(104)} ਦੀ, ਜਿਸ ਨੇ ਹੁਣੇ ਹੁਣੇ ਇਕ
\VUgras{ਕੋਲੇ ਦੀ ਬੋਰੀ} \textsuperscript{(105)} ਪਹੁੰਚਾਉਣ ਨੂੰ ਆਪਣੀ ਗੱਡੀ ਥੋੜੀ ਦੇਰ ਲਈ
ਛੱਡੀ ਹੈ।
\VUsaut{6.0mm}
\VUfilet{22mm}
